(1) 两个势分别满足 \(\phi_{tt}=c_p^2\nabla^2\phi\)、\(\psi_{tt}=c_s^2\nabla^2\psi\)，所以选取衰减分支
\[
p=\sqrt{1-v^2/c_p^2},\qquad q=\sqrt{1-v^2/c_s^2},\qquad 0<q<p.
\]
去除公共 \(e^{i(kx-\omega t)}\)，有 \(u_x=ikAe^{-pkz}+qkBe^{-qkz}\)、\(u_z=-pkAe^{-pkz}+ikBe^{-qkz}\)。由 \(\sigma_{ij}=\lambda(\nabla\cdot\boldsymbol u)\delta_{ij}+2\mu\varepsilon_{ij}\) 得
\[
\sigma_{xz}(0)=-\mu k^2[2ipA+(2-\zeta)B],\quad
\sigma_{zz}(0)=\mu k^2[(2-\zeta)A-2iqB],
\]
其中使用 \(\lambda(p^2-1)+2\mu p^2=\mu(2-\zeta)\)。故非平凡模态须满足
\[
\begin{pmatrix}2ip&2-\zeta\\2-\zeta&-2iq\end{pmatrix}\binom AB=0.
\]
这也说明单独纵波势或横波势无法满足两个零牵引条件，边界局域态必须是二者的混合。

(2) 取矩阵行列式为零，得到严格的 Rayleigh 关系
\[
\boxed{(2-\zeta)^2=4\sqrt{1-\gamma\zeta}\sqrt{1-\zeta}},
\qquad \gamma=\frac{\mu}{\lambda+2\mu}\in(0,\tfrac12].
\]
平方并提取公因子后为 \(\zeta C_\gamma(\zeta)=0\)，其中
\[
C_\gamma(\zeta)=\zeta^3-8\zeta^2+8(3-2\gamma)\zeta-16(1-\gamma).
\]
因为 \(C_\gamma(0)=-16(1-\gamma)<0\)、\(C_\gamma(1)=1>0\)，且
\[
C_\gamma'(\zeta)=3\zeta^2-16\zeta+24-16\gamma
\ge C_\gamma'(1)=11-16\gamma\ge3>0\quad(0\le\zeta\le1),
\]
故存在唯一内部根。该区间原方程两边正数，平方不再引入伪解。端点 \(\zeta=0\) 对应 \(p=q=1\)、\(B=-iA\)，代入位移表达式得到 \(\boldsymbol u\equiv0\)：只是非传播势的规范退化，不能算一支非零表面波。

(3) 令 \(A\) 为实数，边界方程给 \(B=-2ipA/(2-\zeta_R)\)，因此
\[
U_x(z)=ikA\left[e^{-pkz}-\frac{2pq}{2-\zeta_R}e^{-qkz}\right],\quad
U_z(z)=pkA\left[-e^{-pkz}+\frac{2}{2-\zeta_R}e^{-qkz}\right].
\]
显然 \(U_x\) 纯虚而 \(U_z\) 纯实。表面 \(U_z(0)=pkA\zeta_R/(2-\zeta_R)\ne0\)，故实际轨迹为椭圆偏振而不是单一纵横振动。以 \(\rho\ddot u_i=\partial_j\sigma_{ij}\) 乘 \(\dot u_i\)，由应力对称性得到
\[
\partial_t e+\partial_jS_j=0,\quad
e=\frac{\rho}{2}|\dot{\boldsymbol u}|^2+\frac{\lambda}{2}(\nabla\cdot\boldsymbol u)^2+\mu\varepsilon_{ij}\varepsilon_{ij},\quad
S_j=-\sigma_{ij}\dot u_i.
\]
自由面 \(S_z=0\)，无穷深处两种势均指数衰减。对振幅 \(\boldsymbol U(z)\)，令自伴模态弱形式的质量二次型 \(\mathcal M=\int_0^\infty\rho|\boldsymbol U|^2dz\)，弹性二次型 \(\mathcal K_k=\int_0^\infty(\lambda|\nabla_k\cdot\boldsymbol U|^2+2\mu\varepsilon_k^*:\varepsilon_k)dz\)。利用自然自由面边界条件，\(\mathcal K_k=\omega^2\mathcal M\)，对 \(k\) 作驻值求导得
\[
\partial_k\mathcal K_k=2\omega\mathcal M\,d\omega/dk,\quad
\overline{\mathscr E}=\omega^2\mathcal M/2,\quad
\overline{\mathscr S}_x=\omega\,\partial_k\mathcal K_k/4.
\]
所以 \(\overline{\mathscr S}_x/\overline{\mathscr E}=d\omega/dk\)。本问题没有内禀波数尺度，\(\omega=kc_s\sqrt{\zeta_R}\)，故积分平均能量输运速度、群速度与相速都等于 \(c_R\)。这一结论只针对深度积分量，并非任意深度的局域能量密度比。

(4) \(\lambda=0\) 时 \(\gamma=1/2\)，唯一物理根为 \(\zeta_R=0.763932022500210\ldots\)，因此 \(c_R/c_s=0.874032048897642\ldots\)。不可压极限 \(\lambda\to+\infty\) 给 \(\gamma\to0\)，\(\zeta_R=0.912621974615847\ldots\)，\(c_R/c_s=0.955312501025632\ldots\)。两种指数穿透长度为
\[
\ell_p=(k\sqrt{1-\gamma\zeta_R})^{-1},\qquad
\ell_s=(k\sqrt{1-\zeta_R})^{-1},\qquad \ell_s>\ell_p.
\]
均匀无限半空间没有有限厚度的本征长度，故在连续 \(k>0\) 上有线性色散而无几何截止。因 \(c_R<c_s<c_p\)，两种体波势都为倏逝场，不存在向深处传播的体波能流；若误把 \(v\) 取到剪切波速之上，横波势不再指数局域，便失去此束缚波的假设。